%!TEX program = xelatex
\documentclass[aspectratio=169]{beamer}
\usepackage{fontspec}
\usepackage[most]{tcolorbox}

%%%%%%%%%%%% LETRAREN ITURRIA ALDATZEKO (EHUSans) %%%%%%%%%%%%
\setsansfont{EHUSans-Light.otf}[
  Path = {./EHU Sans/},
  BoldFont = EHUSans-Regular.otf,
  ItalicFont = EHUSans-Italic.otf,
  BoldItalicFont = EHUSans-BoldItalic.otf
]

\renewcommand{\familydefault}{\sfdefault}
\newfontface\lightfont{EHUSans-Light.otf}[Path=./EHU Sans/]
\newfontface\blackfont{EHUSans-Black.otf}[Path=./EHU Sans/]
\newfontface\lightitfont{EHUSans-Light.otf}[Path=./EHU Sans/,FakeSlant=0.2]

%%%%%%%%%%%%%%%%%%%%%%%%%%%%%%%%%%%%%%%%%%%%%%%%%%%%%%%%%%%%%%%%%
%%%%%%%%%%%% SIMULACIÓN DE LIGHT ITALIC %%%%%%%%%%%%
\newcommand{\lightit}[1]{{\lightitfont #1}}
%%%%%%%%%%%%%%%%%%%%%%%%%%%%%%%%%%%%%%%%%%%%%%%%%%%%


% Tema clásico académico
\usetheme{CambridgeUS}
\usecolortheme{seahorse}

% Paquetes útiles
\usepackage[spanish]{babel}
\usepackage{graphicx}
\usepackage{amsmath}
\usepackage{siunitx}
\usepackage{booktabs}

% Datos portada
\title[Corrección distorsión óptica]{Desarrollo de un algoritmo para la calibración de la distorsión óptica en adquisición de imágenes satelitales y comprobación experimental}
\author{Xabier Joel Yoldi García}
\institute{Grado en Física \\ EHU}
\date{\today}

\begin{document}

%------------------------------------------------
\begin{frame}
\titlepage
\end{frame}
%------------------------------------------------

\begin{frame}{Índice}
\tableofcontents
\end{frame}

\section{Contexto y motivación}
%-------------------------------------------------
\begin{frame}{Contexto y motivación}
\begin{itemize}
    \item Crecimiento de la observación terrestre mediante microsatélites
    \item Cámaras de alta resolución espacial
    \item Necesidad de precisión geométrica
    \item \lightit{Image Processing}
\end{itemize}
\end{frame}

%-------------------------------------------------
\begin{frame}{Degradaciones en sistemas ópticos}

\begin{columns}[T,onlytextwidth]

\column{0.48\textwidth}

\textbf{Principales fuentes de degradación}

\begin{itemize}
    \item Ruido electrónico
    \begin{itemize}
        \item Lectura
        \item Disparo
        \item Corriente oscura
    \end{itemize}
    \item Aberraciones ópticas
    \begin{itemize}
        \item Esférica, coma, astigmatismo
        \item Curvatura de campo
        \item \alert{\textbf{Distorsión}}
    \end{itemize}
\end{itemize}

\vspace{0.3cm}

\alert{En este trabajo nos centramos en la distorsión geométrica.}

\column{0.52\textwidth}

\textbf{Distorsión óptica}

\vspace{0.3cm}

Deformación geométrica que aumenta con la distancia al centro óptico.

\vspace{0.4cm}

\begin{itemize}
    \item \textbf{Barrel}: líneas se curvan hacia fuera
    \item \textbf{Pincushion}: líneas se curvan hacia dentro
\end{itemize}

\vspace{0.4cm}

Impacto: errores en medida y georreferenciación.

\end{columns}

\end{frame}

\section{Modelado matemático de la distorsión}
%-------------------------------------------------
\begin{frame}{Modelos clásicos de distorsión}

\small
\begin{columns}[T,onlytextwidth]

\column{0.5\textwidth}
\centering
\textbf{Modelo polinómico (Duane)}
\vspace{0.3cm}
\[
R_d = (1 + k_1 R^2 + k_2 R^4 + \dots)\ R
\]


\column{0.5\textwidth}
\centering
\textbf{Modelo de división (Fitzgibbon)}
\vspace{0.25cm}
\[
R = \dfrac{R_d}{1 + \lambda_1 R_d^2 + \lambda_2 R_d^4 + \dots}
\]

\end{columns}
\vspace{0.4cm}
\scriptsize
$R$: posición ideal, sin distorisón\\
$R_d$: posición real, distorsionada
\vspace{0.6cm}

\begin{itemize}
    \item Ambos modelos utilizan expansiones radiales.
    \item Más coeficientes $\Rightarrow$ mayor precisión.
    \item \alert{\textbf{En sistemas ópticos profesionales suele ser suficiente un único coeficiente.}}
\end{itemize}

\end{frame}
%-------------------------------------------------
\begin{frame}{Modelo utilizado}
Modelo radial con un coeficiente:
\begin{equation*}
    R_d = (1 + kR^2)R
\end{equation*}

\begin{itemize}
    \item $R$ : posición ideal
    \item $R_d$ : posición distorsionada
    \item $k$ : coeficiente de distorsión
\end{itemize}

\begin{itemize}
    \item $k > 0$ $\rightarrow$ pincushion
    \item $k < 0$ $\rightarrow$ barrel
\end{itemize}
\end{frame}

%-------------------------------------------------
\begin{frame}{Visualización de la distorsión}
\begin{columns}[T]

\column{0.3\textwidth}
\centering

\includegraphics[width=\linewidth]{Imagenes/imagen_original.jpg}
\\[0.3cm]
\small Imagen ideal

\column{0.3\textwidth}
\centering

\includegraphics[width=\linewidth]{Imagenes/imagen_distorsionada_pincushion.png}
\\[0.3cm]
\small Distorsión pincushion ($k>0$)

\column{0.3\textwidth}
\centering

\includegraphics[width=\linewidth]{Imagenes/imagen_distorsionada_barrel.png}
\\[0.3cm]
\small Distorsión barrel ($k<0$)

\end{columns}

\end{frame}

\section{Método propuesto}
%-------------------------------------------------
\begin{frame}{\lightit{Pipeline} de la metodología}
\begin{columns}[T,onlytextwidth]
  \column{0.52\textwidth}
  \centering
  \textbf{Ejemplo}\\[0.2em]
  \includegraphics[width=\linewidth]{Imagenes/centers.png}\\[-0.2em]


  \column{0.48\textwidth}
  \centering
  \textbf{Pipeline}\\[0.2em]
  \setlength{\tabcolsep}{0.6em}
  \renewcommand{\arraystretch}{1.15}
  \begin{tabular}{c}
  \fbox{\strut Imagen pinhole grid}\\[-0.15em]
  {\Large$\Downarrow$}\\[-0.15em]
  \fbox{\strut Detección de centroides}\\[-0.15em]
  {\Large$\Downarrow$}\\[-0.15em]
  \fbox{\strut Modelado de la placa ideal}\\[-0.15em]
  {\Large$\Downarrow$}\\[-0.15em]
  \fbox{\strut Ajuste por mínimos cuadrados}\\[-0.15em]
  {\Large$\Downarrow$}\\[-0.15em]
  \fcolorbox{black}{gray!10}{\strut \textbf{Estimación del coeficiente $k$}}
  \end{tabular}
\end{columns}
\end{frame}

%-------------------------------------------------
\begin{frame}{Detección de centroides}

\begin{columns}[T,onlytextwidth]

\column{0.48\textwidth}

\textbf{Pipeline de detección}

\begin{enumerate}
    \item Detección de máximos locales
    \item Filtrado por vecindad (eliminación de falsos positivos)
    \item División en ventanas individuales
    \item Umbral adaptativo y binarización
\end{enumerate}

\vspace{0.4cm}

\begin{itemize}
    \item Aislamiento robusto de cada pinhole
    \item Reducción de ruido y artefactos
\end{itemize}

\column{0.52\textwidth}
\centering
\includegraphics[width=0.85\linewidth]{Imagenes/Centroide_candidatos.png}

\end{columns}

\end{frame}

%-------------------------------------------------
\begin{frame}{Cálculo del centroide subpíxel}

\begin{columns}[T,onlytextwidth]

\column{0.45\textwidth}

\small
Centroide como media ponderada de intensidades:

\vspace{0.3cm}

\[
M_i = 
\frac{\sum I(x,y)\, (x,y)}
     {\sum I(x,y)}
\]

\vspace{0.4cm}

\begin{itemize}
    \item Umbral adaptativo por ventana
    \item Eliminación de fondo
    \item \alert{\textbf{Precisión subpíxel}}
\end{itemize}

\column{0.55\textwidth}
\centering
\includegraphics[width=0.8\linewidth]{Imagenes/Centroide_ventana_binaria.png}

\end{columns}

\end{frame}

%-------------------------------------------------
\begin{frame}{Modelado de los pinholes ideales: pipeline}
\centering
\textbf{Flujo de procesado}\\[0.2em]
\setlength{\tabcolsep}{0.6em}
\renewcommand{\arraystretch}{1.15}
\begin{tabular}{@{}c@{}}
\fbox{\strut Placa de pinholes ideal}\\[-0.15em]
{\Large$\Downarrow$}\\[-0.15em]
\fbox{\strut Corrección orientación y desplazamiento}\\[-0.15em]
{\Large$\Downarrow$}\\[-0.15em]
\fbox{\strut Asignación de cada pareja de pinholes}\\[-0.15em]
{\Large$\Downarrow$}\\[-0.15em]
\fcolorbox{black}{gray!10}{\strut \textbf{Parejas de pinholes ideal-real}}
\end{tabular}
\end{frame}

%-------------------------------------------------
\begin{frame}{1) Corrección orientación y desplazamiento}
\begin{columns}[T,onlytextwidth]
    \column{0.5\textwidth}
    \begin{itemize}
        \item Se calculan los ejes principales de la placa utilizando los pinholes cercanos al centro óptico.
        \item Se mantiene fijo el pinhole más cercano al centro óptico.
    \end{itemize}

    \column{0.5\textwidth}
    \centering
    \includegraphics[width=\linewidth]{Imagenes/img_DistMap_0000_real.png}\\
    {\scriptsize Desplazamientos de cada pinhole (flechas zoom x100).}
\end{columns}
\end{frame}

%-------------------------------------------------
\begin{frame}{2) Asignación de cada pareja de pinholes}
\begin{columns}[T,onlytextwidth]
    \column{0.5\textwidth}
    \begin{itemize}
        \item Para la asignación se utiliza el método húngaro de emparejamiento.
        \item Emparejamientos que minimicen la suma de desplazamientos.
        \item A cada pinhole real se le asigna un pinhole ideal. 
    \end{itemize}

    \column{0.5\textwidth}
    \centering
    \includegraphics[width=\linewidth]{Imagenes/img_DistMap_0000_real.png}\\
    \vspace{0.2em}
    {\scriptsize Desplazamientos de cada pinhole (flechas zoom x100).}
\end{columns}
\end{frame}

%-------------------------------------------------
\begin{frame}{Ajuste por mínimos cuadrados}

\begin{columns}[T,onlytextwidth]

\column{0.48\textwidth}

\small
\textbf{Modelo radial}

\vspace{0.2cm}

\[
R_d = (1 + kR^2)R
\]

\vspace{0.4cm}

\textbf{Reordenación}

\[
R_d - R = kR^3
\]

\column{0.52\textwidth}

\small
\textbf{\alert{Linealización}}

\vspace{0.3cm}

\[
y_i = x_i\, k
\]

\vspace{0.3cm}

\[
y = R_d - R
\quad
x = R^3
\]

\vspace{0.6cm}

\centering
\alert{\textbf{Estimación de $k$ mediante ajuste lineal}}

\vspace{0.4cm}

\begin{itemize}
    \item Eliminación de outliers con Z-Score
\end{itemize}

\end{columns}

\end{frame}

\section{Etapa de simulación}
%-------------------------------------------------
\begin{frame}{Entorno de simulación}
\begin{columns}[T,onlytextwidth]
    \column{0.5\textwidth}
    \centering
    \textbf{Flujo de procesado}\\[0.2em]
    \setlength{\tabcolsep}{0.6em}
    \renewcommand{\arraystretch}{1.15}
    \begin{tabular}{@{}c@{}}
    \fbox{\strut Generación de pinhole grid artificial}\\[-0.15em]
    {\Large$\Downarrow$}\\[-0.15em]
    \fbox{\strut Aplicación del método propuesto}\\[-0.15em]
    {\Large$\Downarrow$}\\[-0.15em]
    \fbox{\strut Evaluación de los resultados}
    \end{tabular}
    \vspace{0.4em}
    

    \column{0.5\textwidth}
    \centering
    \includegraphics[width=\linewidth]{Imagenes/img_pinholes.png}
\end{columns}
\end{frame}

%-------------------------------------------------
\begin{frame}{1) Generación de pinhole grid artificial}
\begin{columns}[T,onlytextwidth]
    \column{0.5\textwidth}
    \begin{itemize}
        \item Convolución
        \item PSF aproximada mediante disco de Airy
        \item Aplicación de rotación, desplazamiento y desenfoque
        \item Se aplica el modelo radial con un coeficiente $k$.
        \item Centro óptico configurable para simular desalineaciones.
        \item Remapeo e interpolación para generar la imagen distorsionada.
        \item Adición de ruido para escenarios realistas.
    \end{itemize}

    \column{0.5\textwidth}
    \centering
    \includegraphics[width=0.75\linewidth]{Imagenes/convolucion.png}\\
    \vspace{0.2em}
    {\scriptsize Visualización de convolución.}
\end{columns}
\end{frame}

%-------------------------------------------------
\begin{frame}{2) Aplicación del método propuesto}
\begin{columns}[T,onlytextwidth]
    \column{0.48\textwidth}
    Parámetros variables:
    \begin{itemize}
        \item Coeficiente $k_{real}$
        \item Centro óptico
        \item Rotación
        \item Intensidad
        \item Error en el spacing
    \end{itemize}

    \column{0.52\textwidth}
    \centering
    \begin{tabular}{l c}
    \toprule
    Parámetro variado & Rango estudiado \\
    \midrule
    Coeficiente $k_{real}$ & $[-39.5,\ 39.5]$ m$^{-2}$ \\
    Centro óptico $C_o$ & $\pm 100$ px \\
    Rotación $\theta$ & $\pm 15^\circ$ \\
    Intensidad $I_i$ & $[0,\ 3000]$ \\
    Error en spacing $\delta s$ & 0.01\% \\
    \bottomrule
    \end{tabular}
\end{columns}
\end{frame}

%-------------------------------------------------
\begin{frame}{3) Evaluación de los resultados}
\begin{columns}[T,onlytextwidth]
    \column{0.48\textwidth}
    \includegraphics[width=0.9\linewidth]{Imagenes/RelativeError_vs_k.png}
    \begin{itemize}
        \item Mejor funcionamiento para $|k|>5\text{ m}^{-2}$
    \end{itemize}
    \column{0.52\textwidth}
    \centering
    \begin{tabular}{l c c}
    \toprule
    Escenario & $E_r$ (\%) & $R^2$ \\
    \midrule
    Nominal & 0.28 & 0.9994 \\
    Centro desplazado & 0.70 & 0.9953 \\
    Rotación & 0.32 & 0.9993 \\
    Error spacing (0.01\%) & 7.43 & 0.9964 \\
    Completo + spacing & 9.45 & 0.9882 \\
    \bottomrule
    \end{tabular}
    \vspace{0.5cm}
    \begin{itemize}
        \item Spacing: Condición más sensible
        \item Error del 0.01\% en spacing $\rightarrow$ degradación notable
    \end{itemize}
\end{columns}
\end{frame}

\section{Aplicación en imágenes reales de Satlantis}%-------------------------------------------------
\begin{frame}{Aplicación en las imágenes de Satlantis}
\begin{columns}[T,onlytextwidth]
  \column{0.52\textwidth}
  \centering
  \textbf{Imagen en bruto}\\[0.2em]
  \includegraphics[width=\linewidth]{Imagenes/psfgrid.PNG}\\[-0.2em]


  \column{0.48\textwidth}
  \centering
  \textbf{Pipeline}\\[0.2em]
  \setlength{\tabcolsep}{0.6em}
  \renewcommand{\arraystretch}{1.15}
  \begin{tabular}{c}
  \fbox{\strut Imagen en bruto}\\[-0.15em]
  {\Large$\Downarrow$}\\[-0.15em]
  \fbox{\strut Procesado con Dark Frame}\\[-0.15em]
  {\Large$\Downarrow$}\\[-0.15em]
  \fbox{\strut Determinación del centro óptico}\\[-0.15em]
  {\Large$\Downarrow$}\\[-0.15em]
  \fbox{\strut Aplicación del algoritmo}\\[-0.15em]
  \end{tabular}
\end{columns}
\end{frame}

%-------------------------------------------------
\begin{frame}{1) Procesado con Dark Frame}
\begin{columns}[T]

\column{0.3\textwidth}
\centering
\includegraphics[width=\linewidth]{Imagenes/psfgrid.PNG}
\\[0.3cm]
\small Imagen en bruto

\column{0.3\textwidth}
\centering
\includegraphics[width=\linewidth]{Imagenes/Dark.PNG}
\\[0.3cm]
\small Dark Frame

\column{0.3\textwidth}
\centering
\includegraphics[width=\linewidth]{Imagenes/Pinholes.PNG}
\\[0.3cm]
\small Imagen procesada

\end{columns}

\end{frame}

%-------------------------------------------------
\begin{frame}{2)Determinación del centro óptico}
\begin{columns}[T,onlytextwidth]
    \column{0.48\textwidth}
    \begin{itemize}
        \item Se utiliza el Flat Frame
        \item Se aprovecha el viñeteo y su simetría
        \item Se ajusta un circulo y el centro se determina como el centro óptico.
        \item Otra alternativa si no se dispone del Flat Frame: centro de la imagen
    \end{itemize}

    \column{0.52\textwidth}
    \centering
    \includegraphics[width=\linewidth]{Imagenes/flat_centro_optico.png}
\end{columns}
\end{frame}

%-------------------------------------------------
\begin{frame}{Resultados obtenidos}

\begin{columns}[T,onlytextwidth]
\column{0.48\textwidth}
\centering

\fbox{$\displaystyle k_m = 29.22\ \text{m}^{-2}$}

\vspace{0.6cm}

\fbox{
$\displaystyle
\begin{aligned}
R^2 &= 0.786 \\
\text{RMSE} &= 23.29\ \text{m}^{-2}
\end{aligned}
$
}

\vspace{0.4cm}

\footnotesize
Estimación consistente del coeficiente de distorsión  
en condiciones experimentales reales.
\column{0.52\textwidth}
\centering
\includegraphics[width=\linewidth]{Imagenes/img_DistMap_0000_real.png}

\end{columns}

\end{frame}

\section{Conclusiones}
%-------------------------------------------------
\begin{frame}{Conclusiones}
\begin{itemize}
    \item Desarrollo completo de algoritmo de calibración
    \item Validación exhaustiva mediante simulación
    \item Aplicación en sistema óptico profesional
    \item Alta sensibilidad al spacing
    \item Un coeficiente suficiente en este contexto
\end{itemize}
\end{frame}

%-------------------------------------------------
\begin{frame}{Impacto y trabajo futuro}
\begin{itemize}
    \item Aplicable a calibración industrial
    \item Extensible a modelos con más coeficientes
    \item Posible automatización del proceso
    \item Integración en pipelines de calibración óptica
\end{itemize}
\end{frame}

%-------------------------------------------------
\begin{frame}
\centering
\Huge{Gracias por su atención}
\end{frame}


\end{document}
